\section{Integer bounds for the principal character packet}
\label{sec:sr-character-integers}
The two target estimates of Definition~\ref{def:sr-character-inputs}
concern the same principal host and the same normalized strict lengths.
This section proves their numerical and character-orbit reductions.
The later support arguments make several underlying implications
explicit; the automorphic producers still required are identified at
each use.

\subsection{A single Frobenius prime for the generic character sum}
\label{sec:sr-fixed-seven-generic}
Throughout this subsection the principal host is the field defined in
Section~\ref{sec:sr-arithmetic}, with its fixing subgroup $H_{\rm host}$.
Only its abelian $2$-group character packet, its degree bound, and its
local quadratic characters at the two odd principal primes are used.
We increase the fixed lower threshold for the principal primes to ensure
$q,d>7$. This does not change the order of the uniform quantifiers.

\begin{srproposition}[The exact orbit-integer implication]
\label{prop:sr-orbit-integer-implication}
Fix an imaginary quadratic field $K$. Assume the common trace-integer
construction and the one-character deformation inequality of
Sections~\ref{sec:common-trace-integer}
and~\ref{sec:single-character}, including their stated arithmetic inputs.
There are $c,C\in\mathbb N$, depending only on $K$, with the following
property. For every $a\ge1$, positive modulus $N$, subgroup
$J\le(\mathbb Z/N)^\times$, and prime
$p\nmid30N|\operatorname{disc}K|$, there is a nonzero integer
$Z=Z_{K,J,a,N,p}$ such that
\begin{equation}\label{eq:sr-orbit-integer-size}
 \log|Z|\le c\,\varphi([(\mathbb Z/N)^\times:J])
                    [N(a+1)]^C\log p.
\end{equation}
The integer is chosen before the coefficient prime and coefficient ring.
For every prime $\ell\ne p$ and finite splitting coefficient ring
$\mathcal O$ at $\ell$, if one character of kernel $J$ satisfies the
good-data conditions in Section~\ref{sec:single-character}, all such
characters have finite strict groups and
\begin{equation}\label{eq:sr-orbit-integer-valuations}
       \sum_{\ker\psi=J}v_{\ell,a}(\psi)
                              \le2\operatorname{ord}_\ell Z.
\end{equation}
\end{srproposition}
\begin{proof}
Choose $w\mid p$ and put $Q=Nw$. The complex orbit
$\mathcal O_J$ consists of all characters modulo $N$ with kernel $J$.
Its indexed roots $z_\chi=\chi(Q)^{-1}$ are permuted by every
automorphism of $\mathbb C$, since an injective field homomorphism
preserves a character kernel. If the orbit is nonempty, its quotient
$(\mathbb Z/N)^\times/J$ is cyclic; its faithful characters form one
orbit of cardinality $\varphi([(\mathbb Z/N)^\times:J])$.
All of them are powers $\chi^u$ with $u$ prime to the order.
If the orbit is empty, take $Z=1$.

The trace-integer construction applied to this indexed family gives
a nonzero integer $Z$ satisfying \eqref{eq:sr-orbit-integer-size}.
For every $\ell\ne p$ it controls the product of all the nonzero
differences between admissible traces at inverse Frobenius and the
corresponding Eisenstein traces, for every finite list having pairwise
nonisomorphic members within each orbit position.

Fix such an $\ell$ and $\mathcal O$, and choose an embedding of its
fraction field in $\overline{\mathbb Q}_\ell$ and a field isomorphism
$\iota:\overline{\mathbb Q}_\ell\to\mathbb C$.
Composition preserves kernels and identifies the coefficient-valued
orbit with $\mathcal O_J$. This identification is onto: every character
with a given kernel is a power, prime to the order, of one character
already in the coefficient ring. Good data transfer to every character
of this orbit. This is checked by the same powers: the kernels coincide,
the quadratic local component is preserved, and all the remaining
conditions depend only on the kernel, the order, or the prescribed parity.

Arithmetic Frobenius at $w$ has cyclotomic value $Q$ and Dirichlet
value $\psi(Q)$. Hence its inverse gives the reducible trace
\[
 T_\psi=Q^{-a}\psi(Q)^{-1}+1+Q^a\psi(Q).
\]
The separation supplied by the trace-integer construction permits the
one-character inequality at this inverse Frobenius. It proves finiteness
and supplies a finite set $X_\psi$ of pairwise nonisomorphic admissible
representations with
\[
 v_{\ell,a}(\psi)\le
       2\sum_{\rho\in X_\psi}
          v_\ell\bigl(\operatorname{tr}\rho(\operatorname{Frob}_w^{-1})
                                            -T_\psi\bigr).
\]
Here $v_\ell(\ell)=1$. The product inequality for $Z$ says
\[
 |Z|_\ell\le
       \prod_{\psi}\prod_{\rho\in X_\psi}
       \left|\operatorname{tr}\rho(\operatorname{Frob}_w^{-1})
                                            -T_\psi\right|_\ell .
\]
All factors are nonzero. Taking $-\log_\ell$ bounds the sum of their
valuations by $\operatorname{ord}_\ell Z$. Summing the preceding
one-character inequalities proves
\eqref{eq:sr-orbit-integer-valuations}.
No distinctness is required between different orbit positions: the
integer product retains their separate multiplicities.
The coefficient normalization is exactly
$\operatorname{length}_{\mathcal O}/v_{\mathcal O}(\ell)$, as in
Section~\ref{sec:single-character}.
\end{proof}

\begin{srproposition}[The fixed-seven specialization]
\label{prop:sr-fixed-seven-generic}
Under the arithmetic inputs of
Proposition~\ref{prop:sr-orbit-integer-implication} for the single field
$K=\mathbb Q(\sqrt{-3})$, the generic estimate
\eqref{eq:sr-generic-count} holds with one natural exponent fixed
before the host, twist range, coefficient rings and finite prime set.
More precisely, if $c,C$ are the constants above, choose
$k\in\mathbb N$ once and for all with
\[
                2^k\ge\max(1,2c\log7).
\]
Then one may use $A=C+1+k$.
\end{srproposition}
\begin{proof}
Fix $t\ge2$, a principal host with $q,d>7$, and $1\le a<t$.
Let $\mathcal J_a$ be the set of kernels of its complex characters
in $\mathcal X_a$. The entire kernel-$J$ orbit is contained in
$\mathcal X_a$ for every $J\in\mathcal J_a$.
Indeed, a character of the same kernel is $\psi^u$ for $u$ prime to
$r=\operatorname{ord}\psi$. It is trivial on $H_{\rm host}$.
If $r>1$ it is a power of two, so $u$ is odd and
$\psi^u(-1)=\psi(-1)$; if $r=1$ the orbit is a singleton.
Consequently these orbits partition $\mathcal X_a$, and
\begin{equation}\label{eq:sr-orbit-totient-sum}
 \sum_{J\in\mathcal J_a}
           \varphi([(\mathbb Z/N)^\times:J])
             =|\mathcal X_a|\le[E:\mathbb Q]\le N.
\end{equation}
This equality concerns character kernels in the full group of units
modulo $N$, with their common host fixing subgroup retained.

Put $p=7$. The choice $q,d>7$ gives $7\nmid N$, while
$|\operatorname{disc}K|=3$, so $7\nmid30N|\operatorname{disc}K|$.
Proposition~\ref{prop:sr-three-auxiliary} verifies the good-data
conditions for every odd nonresonant coefficient prime and every
character in $\mathcal X_a$. For the source's character $\psi'$ with
the $\ell$-local component removed, the final kernel condition follows
from inertia at three as well: its conductor still divides $N$.
Use Proposition~\ref{prop:sr-orbit-integer-implication} once for
each $J\in\mathcal J_a$, obtaining $Z_J=Z_{K,J,a,N,7}$.
Every $\ell$ in the prescribed generic set $S$ differs from seven.
Thus all the following strict lengths are finite and
\begin{align*}
 \sum_{\ell\in S}\log\ell
                   \sum_{\psi\in\mathcal X_a}v_{\ell,a}(\psi)
 &\le2\sum_{J\in\mathcal J_a}
                     \sum_{\ell\in S}
                      (\log\ell)\operatorname{ord}_\ell Z_J\\
 &\le2\sum_{J\in\mathcal J_a}\log|Z_J|\\
 &\le2c\log7\,N[N(a+1)]^C .
\end{align*}
The middle inequality is prime factorization of the nonzero integers
$Z_J$; it does not contain $|S|$.
Put $M=Nt$. Since $a<t$, we have $N(a+1)\le M$, $N\le M$,
and $M\ge2$. Our fixed choice of $k$ gives
\[
        2c\log7\,N[N(a+1)]^C
                         \le M^{C+1+k}.
\]
This proves the asserted uniform quantifier order.
\end{proof}

\paragraph{Source verification and remaining scope.}
The exact proposition $\srcid{PropGK}$ is defined in
\href{https://github.com/ChristianKudeba/alpoge-rh-verification/blob/alpoge-source-7cbf7a0c/src/Development/AutomorphicArithmetic/AutomorphicVarieties/UnitaryThreeGroups/KatoUse.lean}
{\srcid{UnitaryThreeGroups.KatoUse}}.
Its producer is
\href{https://github.com/ChristianKudeba/alpoge-rh-verification/blob/alpoge-source-7cbf7a0c/src/Development/AutomorphicArithmetic/AutomorphicVarieties/ModularityArgument/ThresholdArguments/DenseThresholdCases/FPropGKU.lean}
{\srcid{FPropGKU.propGK_of_numberU3U_of_oneCharU}}.
The two arguments above prove the orbit assembly and the entire
coefficient-prime summation. Their stated arithmetic input still
includes the unramified automorphy used by the one-character inequality.
In the archive its producer runs through
\srcid{DnEndsTop.ends0_alt}, \srcid{LinkedEndE},
\srcid{PrLinkedEnd}, \srcid{AbsMain}, and the characteristic-$\ell$
component-propagation theorem \srcid{PropagationTheoremGL3}.
The fixed-seven specialization does not replace that producer.

\subsection{Arithmetic absorption in the principal exceptional packet}
\label{sec:sr-exceptional-absorption}

Throughout this subsection the principal host is the actual fixed field
used in Section~\ref{sec:sr-arithmetic}:
\[
 E=\mathbb Q\bigl(\zeta_{2^{h+2}}+\zeta_{2^{h+2}}^{-1},
                   \sqrt{-q},\sqrt{-d}\bigr),\qquad N=2^{h+2}qd.
\]
Equivalently, its fixing subgroup at the dyadic factor is generated by
\(-1\). We use precisely its character packet \(\mathcal X_a\).
In particular the dyadic component of an admitted character is even.
Enlarging this field to the full dyadic cyclotomic compositum would
enlarge that packet and is not a step in the following argument.
The primes \(q,d\) satisfy the principal congruences, exceed \(2t+1\),
and exceed any fixed thresholds stated below.

For clarity, the assertions below concerning a deformation comparison
are implications from explicitly stated comparisons. They establish the
uniform scalar absorptions, not the construction of the deformation
points or their automorphic realization.

\begin{srproposition}\label{prop:sr-principal-case-II-arithmetic}
Let \(1\leq a<t\), let \(\psi\in\mathcal X_a\) be nontrivial, and suppose
\[
 \ell\in\{3,7\}\quad\hbox{or}\quad
                    (a\geq2,\ \ell-1\mid2a).
\]
Assume \(q,d>7\). Then \(\ell\nmid N\), the order of \(\psi\) is a
nontrivial power of two, and its local components at odd primes have
order at most two. Consequently, when \(a\geq2\), this is Case II
of the source; neither Case III nor Case IV is needed for this packet.

Write \(c=c_2(\psi)\) for the exponent in the conductor \(2^c\) of its
dyadic component. Remove that component from \(\psi\), and call the
order of the remaining character \(b\). Then \(b\mid2\), and
\[
 f=\operatorname{ord}_{\operatorname{lcm}(4,b)}(\ell)
       =\operatorname{ord}_4(\ell)\leq2,\qquad
 w=v_2(\ell^f-1),\qquad 4^w<\ell^4 .
\]
\end{srproposition}
\begin{proof}
The odd ramification of \(E/\mathbb Q\) comes from the two quadratic
radicals, so the components at \(q,d\) have exponent two; there are no
other odd ramified components. The dyadic quotient is a \(2\)-group.
Thus every host character has \(2\)-power order, and its order is even
when it is nontrivial.

For an exceptional prime, \(\ell<q,d\) follows from \(q,d>7\).
For a resonant prime, \(\ell\leq2a+1<q,d\).
Since \(\ell\) is odd, in either case \(\ell\nmid2^{h+2}qd=N\).
The \(\ell\)-power-order part of \(\psi\) and its local component at
\(\ell\) are therefore both trivial. Hence the character denoted
\(\mathrm{chiP}\) in the source is exactly \(\psi\), and its order is
even. Together with \(a\geq2\) and the displayed exceptional or resonant
alternative, these are exactly the three conditions defining Case II.
Case III requires a nontrivial character of odd order, which cannot
occur here. Case IV also requires odd order after taking \(\mathrm{chiP}\).

Removing the dyadic component leaves the product of the two quadratic
components, whose order divides two. Thus the modulus defining \(f\)
is four. Finally \(2^w\leq\ell^f-1<\ell^2\), which gives \(4^w<\ell^4\).
\end{proof}

\begin{srproposition}[The small-conductor twist-one branch]
\label{prop:sr-small-conductor-exceptional}
At \(a=1\), the threshold
\[
                    R_1(\ell,\psi)=16\ell^2\,4^w
\]
equals \(9216\) for \(\ell=3\) and \(200704\) for \(\ell=7\).
If \(2^{c_2(\psi)}<R_1(\ell,\psi)\), then
\[
 \operatorname{ord}\psi\leq
 \begin{cases}2^{11},&\ell=3,\\2^{15},&\ell=7.\end{cases}
\]
Consequently this entire branch satisfies
\[
 (\log\ell)v_{\ell,1}(\psi)\leq(\ell N)^{B_{\mathrm{small}}},
 \qquad
 B_{\mathrm{small}}=196609\cdot393216^{24}.
\]
The exponent is fixed before the host, character and splitting ring.
\end{srproposition}
\begin{proof}
Both three and seven are \(3\bmod4\), so \(f=2\).
The valuations \(v_2(3^2-1)=3\) and \(v_2(7^2-1)=4\) give the displayed
thresholds. Since
\[
 2^{13}<9216<2^{14},\qquad
 2^{17}<200704<2^{18},
\]
the strict threshold bounds \(c_2(\psi)\) by thirteen or seventeen.

For \(c\geq3\), the exponent of \((\mathbb Z/2^c)^\times\) is
\(2^{c-2}\). Indeed the units are generated by \(-1\) and \(5\), and
\(v_2(5^{2^j}-1)=j+2\), by induction using the factorization
\(x^2-1=(x-1)(x+1)\). For \(c=0,1,2\), the unit-group exponent is at
most two. The odd local components also have exponent at most two.
The order of their product is the least common multiple of these
\(2\)-power orders, proving the two order bounds.

Apply Proposition~\ref{prop:sr-bounded-order-one} with \(R=2^{15}=32768\).
Its exponent \((6R+1)(12R)^{24}\) is the displayed integer.
This also proves finiteness in this branch.
\end{proof}

\paragraph{Normalization under a change of splitting ring.}
If \(\mathcal O\subset\mathcal O'\) is a finite extension of splitting
discrete valuation rings and \(\psi'\) is the same character in
\(\mathcal O'\), then
\[
                         v_{\ell,a}(\psi')=v_{\ell,a}(\psi).
\]
Indeed \(\mathcal O'\) is finite free over \(\mathcal O\), and
\[
 (Q/\mathcal O)\otimes_{\mathcal O}\mathcal O'
                       \simeq Q'/\mathcal O'.
\]
Choosing an \(\mathcal O\)-basis identifies the continuous cochain
complex, and every local restriction complex, with a finite direct
sum of the old complex. Thus cohomology and the simultaneous local
kernel commute with this scalar extension. If
\(e'=v_{\mathcal O'}(\varpi_{\mathcal O})\), the cyclic decomposition
of a finite torsion module \(T\) gives
\[
 \operatorname{length}_{\mathcal O'}(T\otimes\mathcal O')
       =e'\operatorname{length}_{\mathcal O}T,\qquad
 v_{\mathcal O'}(\ell)=e'v_{\mathcal O}(\ell).
\]
Dividing gives the asserted identity. Finiteness descends by faithful
flatness as well as ascending by scalar extension. In particular,
adjoining the auxiliary dyadic roots of unity creates no factor in
the normalized estimate.

\begin{srproposition}[Uniform absorption of the rank-four comparison]
\label{prop:sr-rank-four-absorption}
Fix natural constants \(C,A\) in the rank-four common-integer
comparison, before the principal host is chosen.
For each Case II instance with \(a\geq2\), set
\[
 u=\max\{\ell N4^w,\ 2^{c_2(\psi)+6}\ell^3\},\quad
 n=\lceil\log_2u\rceil,\quad
 M=\ell\,2^nN,\quad k=2(\ell-1)2^n.
\]
Suppose the following comparisons have been constructed for these
same data:
\begin{enumerate}
\item At an eligible prime \(p\geq5\), a finite pairwise nonisomorphic
family of admissible rank-four representations \(\mathcal R\), with
nonzero trace differences \(d_\rho\), satisfies
\[
 (\log\ell)v_{\ell,a}(\psi)
                  \leq2\sum_{\rho\in\mathcal R}-\log|d_\rho|_\ell .
\]
\item A nonzero integer \(J\), for the common root-of-unity order \(k\),
satisfies
\[
 |J|_\ell\leq\prod_{\rho\in\mathcal R}|d_\rho|_\ell,\qquad
 \log|J|\leq C\varphi(k)\{M(a+1)\}^{A}\log p.
\]
\end{enumerate}
These are the normalizations in
\(\mathrm{Raw.Ineq}\), \(\mathrm{NumberSp4Alt}\), and
\(\mathrm{CountS.exists\_frob\_nonvanishing\_with\_sha\_bound}\).
There is no additional factor \(a+1\) outside the indicated \(A\)-th
power.

Choose once and for all \(b_4\in\mathbb N\) with \(2C\leq2^{b_4}\).
If the two comparisons hold at every eligible prime, they imply
\[
 (\log\ell)v_{\ell,a}(\psi)\leq(\ell Na)^{B_4},\qquad
                         B_4=10A+12+b_4 .
\]
\end{srproposition}
\begin{proof}
Put \(x=\ell Na\). The conductor \(2^{c_2(\psi)}\) divides \(N\).
By Proposition~\ref{prop:sr-principal-case-II-arithmetic},
\[
 u\leq\max\{\ell^5N,\ 64N\ell^3\}\leq\ell^7N.
\]
Here \(64\leq\ell^4\) since \(\ell\geq3\).
For every positive \(u\), \(2^{\lceil\log_2u\rceil}\leq2u\);
hence
\[
                    2^n\leq2\ell^7N\leq\ell^8N.
\]
It follows that
\[
 \begin{split}
 k&=2(\ell-1)2^n\leq2\ell^9N\leq x^{10},\\
 M(a+1)&\leq\ell^9N^2(2a)
                    \leq\ell^{10}N^2a\leq x^{10}.
 \end{split}
\]

Choose a prime divisor \(p\) of the positive integer \(6\ell N-1\).
It is different from two and three, so \(p\geq5\), and
\(p\nmid\ell N\). Since the only added factor in \(M\) is a power
of two, \(p\nmid M\). It is also outside the Case II ramification
set, whose prime divisors are those of \(2\ell N\).
Moreover
\[
       p\leq6\ell N-1<6\ell N\leq x^2,\qquad \log p\leq x^2,
\]
where \(a\geq2\), \(\ell\geq3\) ensure \(\ell Na^2\geq6\).

Take negative logarithms of the product comparison. Since the
differences are nonzero,
\[
 (\log\ell)v_{\ell,a}(\psi)
       \leq2v_\ell(J)\log\ell
       \leq2\log|J|
       \leq2Ck\{M(a+1)\}^{A}\log p
       \leq2C x^{10A+12}.
\]
The integer inequality \(v_\ell(J)\log\ell\leq\log|J|\) is immediate
from \(\ell^{v_\ell(J)}\mid J\). Finally
\(2C\leq2^{b_4}\leq x^{b_4}\), proving the stated exponent.
All constants precede the host, \(\ell,a,\psi\), and coefficient
extension.
\end{proof}

\begin{srproposition}[A fixed Gaussian test prime]
\label{prop:sr-gaussian-fixed-prime}
Raise the fixed principal-host threshold to \(q,d>11\).
Then \(p=11\) is one eligible inert test prime for every Gaussian
twist-one instance at \(\ell=3,7\). Its unique Gaussian prime \(w\)
has norm \(121\), and \(11\nmid\ell N\).

Suppose the Gaussian common-integer comparison has fixed constants
\(C_G,A_G\in\mathbb N\), so that, with the Galois-stable list of
all complex characters modulo \(N\), it supplies \(J_G\ne0\) with
\[
              \log|J_G|
                 \leq C_G\varphi(N)(2N)^{A_G}\log11 .
\]
Suppose also that the marked universal ring, its finite weight-zero
part, the automorphy and rigidity of its type-I points, and the
cotangent comparison have supplied the source's inequality
\[
 v_{\ell,1}(\psi)
     \leq2\sum_{\rho} -\log_\ell|d_\rho|_\ell ,
\]
with the nonisomorphism conditions needed for the common-integer
comparison after restriction to \(G_{\mathbb Q(i)}\).
Then
\[
 (\log\ell)v_{\ell,1}(\psi)
        \leq4C_G\varphi(N)(2N)^{A_G}\log11
        \leq(\ell N)^{B_G},
 \qquad B_G=A_G+1+b_G,
\]
where \(b_G\) is any fixed natural number satisfying
\(4C_G2^{A_G}\log11\leq3^{b_G}\).
\end{srproposition}
\begin{proof}
Since \(11\equiv3\bmod4\), the polynomial \(X^2+1\) is irreducible
over \(\mathbb F_{11}\): a root would give an element of order four
in the group \(\mathbb F_{11}^{\times}\), whose order is ten.
Thus \(11\) is inert in \(\mathbb Q(i)\) and its residue-field norm
is \(11^2\). It divides neither the discriminant four nor \(N\);
also \(\ell\ne11\).

A point's automorphic realization may choose either of the two
embeddings \(\mathbb Q(i)\hookrightarrow\Omega\). Partition the
finite representative family accordingly. For each part the
common-integer comparison bounds the sum of the negative
\(\ell\)-adic logarithms by \(v_\ell(J_G)\).
The same integer works for either embedding because its choice
precedes the algebraic closure, embedding and coefficient prime
in that comparison. Adding the two inequalities and then applying
the cotangent factor two gives
\[
             (\log\ell)v_{\ell,1}(\psi)
                         \leq4v_\ell(J_G)\log\ell
                         \leq4\log|J_G|.
\]
This argument uses the source's distinctness after restriction;
pairwise distinct rational representations alone would not
justify this step.

Finally \(\varphi(N)\leq N\), so the right side is at most
\(4C_G2^{A_G}\log11\cdot N^{A_G+1}\).
Since \(\ell\geq3\), \(N\geq1\), and \(b_G\geq0\), this is at most
\((\ell N)^{A_G+1+b_G}\).
\end{proof}

Combining the preceding three branches would permit the single
exponent
\[
                  B=\max\{B_{\mathrm{small}},B_4,B_G\}.
\]
The small-conductor twist-one branch is proved above without
deformation inputs. The other two displayed absorptions require
their stipulated deformation and common-integer comparisons.
They do not assert that the finite weight-zero, no-block,
automorphic-point or marked-ring constructions have been supplied.

\paragraph{Frozen-source locations.}
The precise source definitions are
\(\mathrm{HeightAndVanishing.DirichletPrimeToLPartMulLPowerPart}\)
(\(\mathrm{Host.H}\), \(\mathrm{R1}\), \(\mathrm{nII}\), and the cases),
\(\mathrm{DenseThresholdCases.AfStmt}\)
(\(\mathrm{NumberSp4Alt}\)),
\(\mathrm{DenseThresholdCases.CountS}\),
\(\mathrm{DenseThresholdCases.ConductorAtMapInjectiveCoeff}\),
\(\mathrm{DenseThresholdCases.EHostNumberU}\), and
\(\mathrm{DenseThresholdCases.CountVOf}\), at tag
\texttt{alpoge-source-7cbf7a0c}.
The finite weight-zero/no-block/automorphic rank-four producer is
\(\mathrm{Sp4In0.sp4In0\_alt}\), through
\(\mathrm{CuFinal.main4\_alt}\).
The Gaussian marked-ring producer is
\(\mathrm{DnEndsTop.ends0\_alt}\), second component
\(\mathrm{ThmS6Core\_s}\).
